\documentclass[12pt,a4paper]{article}
\usepackage[utf8]{inputenc}
\usepackage[T1]{fontenc}
\usepackage{amsmath,amssymb,amsfonts,mathtools}
\usepackage{geometry}
\geometry{left=2.5cm,right=2.5cm,top=2.5cm,bottom=2.5cm}
\usepackage{booktabs}
\usepackage{tabularx}
\usepackage{array}
\usepackage{hyperref}
\usepackage{url}
\usepackage{graphicx}
\usepackage{xcolor}
\usepackage{titlesec}
\usepackage{fancyhdr}
\usepackage{microtype}
\usepackage{fontspec}
\usepackage{parskip}
\usepackage{float}
\usepackage{physics}
\usepackage{bm}
\usepackage{xeCJK}  % 中文支持

% 中文字体（根据系统调整）
\setCJKmainfont{SimSun}
\setCJKsansfont{SimSun}
\setCJKmonofont{SimSun}
\setmainfont{Times New Roman}

% YUCT macros
\newcommand{\Keff}{K_{\mathrm{eff}}}
\newcommand{\kappac}{\kappa_c}
\newcommand{\eps}{\varepsilon}
\newcommand{\betaf}{\beta}
\newcommand{\dint}{D\!+\!I\!\cdot\!R}
\newcommand{\Sodd}{S_{\mathrm{odd}}}
\newcommand{\Seven}{S_{\mathrm{even}}}

\titleformat{\section}{\Large\bfseries}{\thesection}{1em}{}
\titleformat{\subsection}{\large\bfseries}{\thesubsection}{1em}{}

\hypersetup{
	colorlinks=true,
	linkcolor=blue,
	citecolor=blue,
	urlcolor=blue,
	pdftitle={YUCT 与耗散结构：基于协调的观点},
	pdfauthor={Alexey V. Yakushev}
}

\begin{document}
	
	\title{\huge\textbf{YUCT 与耗散结构}\\[0.2cm]
		\Large 非平衡失稳的基于协调的解释\\[0.3cm]
		\large \textit{连接普适误差定律与对流失稳}}
	\author{Alexey V. Yakushev\\[0.2cm] \small Yakushev Research \quad \url{https://yuct.org/}}
	\date{\small 2026年5月 \quad 版本 1.0}
	\maketitle
	
	\begin{abstract}
		\noindent 
		雅库舍夫统一协调理论 (YUCT) 通过单一参数——协调效率 \(K_{\mathrm{eff}}\) 以及普适误差标度律 \(\varepsilon = \kappa_c\alpha K_{\mathrm{eff}}^{-\beta}\)（\(\beta = 2/3\)，\(\kappa_c = 1/3\)）来描述有序复杂性。本附录探讨该定律对非平衡定态稳定性（特别是 B\'{e}nard 对流的起始）的推论。我们提出一种将简单流体的 \(K_{\mathrm{eff}}\) 用其热导率和剪切粘度定义的方案——这些量均可通过常规实验测量。利用该定义，我们将熵产生率 \(\sigma\) 表达为 \(K_{\mathrm{eff}}\) 的函数，并表明线性稳定性阈值对应于协调效率的一个普适临界值 \(K_{\mathrm{eff}}^{\mathrm{crit}} = (S_{\mathrm{odd}}/S_{\mathrm{even}})^{1/\beta} \approx 1.837\)。该值直接源于 YUCT 的代数循环（\(S_{\mathrm{odd}} = 6/5,\; S_{\mathrm{even}} = 4/5,\; \beta = 2/3\)），不包含任何自由参数。预测表明：对于任何牛顿流体，当按本文定义的其有效协调效率低于约 \(1.837\) 时，对流将开始发生。我们推导出一个关于临界 Rayleigh 数与流体 Prandtl 数之间的定量、可证伪的关系，这一关系可通过现有实验数据进行检验。推导过程没有使用 KPZ 关系或分形维数，仅依赖于 YUCT 的代数结构。本文明确讨论了当前唯象处理的局限性。
	\end{abstract}
	\vspace{0.5cm}
	\noindent\textbf{关键词:} YUCT，B\'{e}nard 对流，协调效率，熵产生，Prandtl 数，代数循环。
	
	\newpage
	\tableofcontents
	\newpage
	
	\section{引言}
	
	雅库舍夫统一协调理论 (YUCT) 基于这样一个前提：协调——系统组件利用预先共享的字典对齐其状态的能力——是物理定律涌现的基本过程~\cite{Yakushev2026}。其预测能力建立在一个关于协调相对误差的单一标度律上：
	\begin{equation}\label{eq:universal}
		\varepsilon = \kappa_c\,\alpha\, K_{\mathrm{eff}}^{-\beta},\qquad 
		\beta = \frac{2}{3},\quad \kappa_c = \frac{1}{3},
	\end{equation}
	其中 \(\alpha\) 是一个约为 1 的系统特定常数。\(\beta\) 和 \(\kappa_c\) 的值不是自由参数；它们来自分层协调网络的代数自洽性，详见附录~\cite{AppendixY, AppendixL, AppendixFerra, AppendixAF}。由此产生的代数循环给出两个精确的有理常数：
	\begin{equation}\label{eq:S-constants}
		S_{\mathrm{odd}} = \frac{6}{5} = 1.2,\qquad
		S_{\mathrm{even}} = \frac{4}{5} = 0.8,
	\end{equation}
	它们满足 \(S_{\mathrm{odd}} + S_{\mathrm{even}} = 2\) 且 \(S_{\mathrm{odd}}/S_{\mathrm{even}} = 3/2 = 1/\beta\)。这些数字是所有后续推导的基石。
	
	在另一条研究线中，Ilya Prigogine 提出的耗散结构理论描述了远离热力学平衡的开放系统如何自发演化到更高复杂性的状态~\cite{Prigogine1977, Glansdorff1971}。其典型例子是 B\'{e}nard 对流：从下方加热的水平流体层在温度梯度超过临界值时，会发展出规则的滚动图案~\cite{Chandrasekhar1961}。Prigogine 的稳定性判据指出，当过量熵产生 \(\delta_X P\) 变为正时，定态变得不稳定。
	
	本附录的目的是表明 Prigogine 的判据可以用 YUCT 的语言重新表述，并且 YUCT 的代数循环对对流起始给出了一个定量的、可检验的预测。该推导是唯象的，但除了 YUCT 框架本身已有的参数外，不包含任何自由参数。我们也明确承认当前方法的局限性，并概述了建立完整微观理论所需的步骤。
	
	\section{简单流体的协调效率}
	
	要将 YUCT 应用于物理系统，必须识别出相关的“字典”和“索引”，并由此计算 \(K_{\mathrm{eff}}\)。对于处于温度梯度中的流体，热输运可视为一系列局域协调事件：较高温度的分子通过碰撞将能量传递给较低温度的分子。该过程的效率受到两个因素的限制：介质传导热的能力（热导率 \(\lambda\)）和耗散能量的内摩擦（剪切粘度 \(\eta\)）。具有高 \(\lambda\) 和低 \(\eta\) 的流体是“好的协调器”，因为能量被快速移动且损失很小。
	
	因此我们提出以下可操作的定义：
	\begin{equation}\label{eq:Keff-fluid}
		K_{\mathrm{eff}} \equiv \frac{\lambda}{\lambda_0} \left( \frac{\eta_0}{\eta} \right)^{\gamma},
	\end{equation}
	其中 \(\lambda_0\) 和 \(\eta_0\) 是参考值（例如 \(20^\circ\)C 下的水），指数 \(\gamma\) 被选择使得 \(K_{\mathrm{eff}}\) 成为一个无量纲、尺度无关的量。由于 \(\lambda/\eta\) 的量纲为 \((\mathrm{长度})^2/(\mathrm{时间}\cdot\mathrm{温度})\)，一个自然的选择是 \(\gamma = 1\)，于是：
	\begin{equation}\label{eq:Keff-fluid2}
		K_{\mathrm{eff}} = \frac{\lambda / \lambda_0}{\eta / \eta_0}.
	\end{equation}
	这个定义类似于附录~C 中为化学元素引入的“配位数”，它使 \(K_{\mathrm{eff}}\) 可以直接测量。
	
	\section{熵产生与普适误差定律}
	
	对于静止流体中的温度梯度 \(\nabla T\)，唯一的热力学通量是热流 \(J = -\lambda \nabla T\)，共轭力为 \(X = \nabla(1/T)\)。单位体积的熵产生率为：
	\begin{equation}\label{eq:sigma}
		\sigma = J \cdot \nabla\left(\frac{1}{T}\right) \approx \frac{\lambda |\nabla T|^2}{T^2}.
	\end{equation}
	在 YUCT 的图像中，热流是许多微观协调行为的结果。其中一部分 \(\varepsilon\) 未能相干地传递能量，因此有效通量减少：
	\begin{equation}\label{eq:J-effective}
		J_{\mathrm{eff}} = J \; \bigl(1 - \varepsilon\bigr).
	\end{equation}
	“损失”的通量对熵产生有贡献，但对梯度的相干移除没有贡献。利用普适误差定律 \eqref{eq:universal}，我们可以写出：
	\begin{equation}\label{eq:sigma-K}
		\sigma = \frac{\lambda |\nabla T|^2}{T^2} \frac{1}{1 - \kappa_c\alpha K_{\mathrm{eff}}^{-\beta}}.
	\end{equation}
	当 \(K_{\mathrm{eff}} \gg 1\) 时，修正项可忽略，系统遵守通常的线性傅里叶定律。随着 \(K_{\mathrm{eff}}\) 减小（流体相对于其导热性变得更粘稠），误差增大，对于相同的施加梯度，熵产生增加。
	
	\section{稳定性阈值与临界协调效率}
	
	Prigogine 的稳定性判据指出，当涨落导致过量熵产生变为正值时，定态变得不稳定~\cite{Glansdorff1971}。在 YUCT 表述中，涨落可视为流体协调状态的局部变化——即 \(K_{\mathrm{eff}}\) 对其定态值的临时偏离。利用 \eqref{eq:sigma-K}，小的变化 \(\delta K_{\mathrm{eff}}\) 导致熵产生的变化：
	\begin{equation}
		\delta_X P \propto - \delta K_{\mathrm{eff}}.
	\end{equation}
	因此，使 \(K_{\mathrm{eff}}\) \emph{减小} 的涨落会使 \(\delta_X P > 0\) 并倾向于被放大。这种放大加速了协调的退化，驱动系统走向新的状态。
	
	因此，当 \(K_{\mathrm{eff}}\) 低于临界值时，静止导热态变得不稳定。由于 YUCT 代数循环固定了所有无量纲比值，该临界值必须能用基本常数表示：
	\begin{equation}\label{eq:Kcrit}
		K_{\mathrm{eff}}^{\mathrm{crit}} = \left( \frac{S_{\mathrm{odd}}}{S_{\mathrm{even}}} \right)^{1/\beta}
		= \left( \frac{3}{2} \right)^{3/2} \approx 1.837.
	\end{equation}
	指数 \(1/\beta\) 源于这样的要求：阈值处的误差参数 \(\varepsilon\) 应为 1 的量级（系统不再能区分有用的协调与噪声）。比值 \(S_{\mathrm{odd}}/S_{\mathrm{even}} = 3/2\) 是 YUCT 的普适有序‑无序比。方程 \eqref{eq:Kcrit} 不包含自由参数。
	
	\section{对临界 Rayleigh 数的预测}
	
	深度为 \(d\)、密度为 \(\rho\)、比热为 \(c_p\)、热膨胀系数为 \(\alpha_T\)、施加温差为 \(\Delta T\) 的流体层的 Rayleigh 数为：
	\begin{equation}
		\mathrm{Ra} = \frac{\rho^2 c_p \alpha_T g \Delta T d^3}{\lambda \eta}.
	\end{equation}
	用我们的协调效率 \eqref{eq:Keff-fluid2}，可以写成：
	\begin{equation}
		\mathrm{Ra} = \frac{A}{K_{\mathrm{eff}}},
	\end{equation}
	其中 \(A\) 是一个仅依赖于几何和边界条件、而不依赖于流体性质的常数。通过此识别，阈值条件 \(K_{\mathrm{eff}} = K_{\mathrm{eff}}^{\mathrm{crit}}\) 转化为临界 Rayleigh 数：
	\begin{equation}\label{eq:Racrit}
		\mathrm{Ra_c} = \frac{A}{K_{\mathrm{eff}}^{\mathrm{crit}}}.
	\end{equation}
	如果能通过一种参考流体（例如水）确定常数 \(A\)（水的经典临界 Rayleigh 数对于刚性边界约为 \(\mathrm{Ra}_{\mathrm{水}} \approx 1708\)，而我们的定义给出某个 \(K_{\mathrm{eff}}^{\mathrm{水}}\)），那么任何其他流体的临界 Rayleigh 数都可以预测：
	\begin{equation}\label{eq:Racrit-predict}
		\boxed{ \mathrm{Ra_c}(\text{流体}) = \mathrm{Ra_c}(\text{水}) \;
			\frac{K_{\mathrm{eff}}^{\mathrm{水}}}{K_{\mathrm{eff}}(\text{流体})} }.
	\end{equation}
	利用 \eqref{eq:Keff-fluid2}，这变为临界 Rayleigh 数与流体热导率和粘度之间的关系：
	\begin{equation}
		\mathrm{Ra_c} \propto \frac{\eta / \eta_0}{\lambda / \lambda_0}.
	\end{equation}
	因此，具有更高粘度（协调性较差）的流体应该具有 \emph{更大} 的临界 Rayleigh 数——使其更难产生对流。该预测在定性上与经典结果一致：粘稠的油类具有比水高得多的临界 Rayleigh 数。定量检验需要测量流体的 \(\lambda\) 和 \(\eta\)，根据 \eqref{eq:Keff-fluid2} 计算 \(K_{\mathrm{eff}}\)，并将预测的 \(\mathrm{Ra_c}\) 与实验观察到的对流起始值进行比较。在单一参考流体标定之后，不再有可调参数。
	
	\section{可检验的推论与实验方案}
	
	\subsection{用现有数据进行的即时检验}
	文献包含各种流体的临界 Rayleigh 数的大量表格（例如参见~\cite{Chandrasekhar1961, Cross1993} 的综述）。利用相关温度下测量的 \(\lambda\) 和 \(\eta\) 值，可以计算每种流体的 \(K_{\mathrm{eff}}\)，并验证比值 \(\mathrm{Ra_c}/\mathrm{Ra_c}^{\mathrm{水}}\) 是否等于 \(K_{\mathrm{eff}}^{\mathrm{水}} / K_{\mathrm{eff}}\)。若 \(R^2 > 0.9\) 的正相关将为 YUCT 机制提供强有力的证据。
	
	\subsection{针对性实验}
	可以对两种协调效率差异很大的流体（例如水和硅油）进行受控实验。对于每种流体，高精度测量临界温差 \(\Delta T_c\)。YUCT 的预测是：在固定几何条件下，临界温差之比应等于由 \eqref{eq:Keff-fluid2} 定义的协调效率之反比。这是一个无参数、可证伪的预测。
	
	\section{局限性与未解决的问题}
	
	当前的分析是唯象的。完整的微观理论需要从第一原理推导出协调误差 \(\varepsilon\) 与热力学通量之间的联系，而不是假设 \eqref{eq:J-effective}。此外，\eqref{eq:Keff-fluid2} 中对 \(K_{\mathrm{eff}}\) 的定义是基于动机的设想；其他定义也可能成立，应与数据进行比较。分析还忽略了化学反应的作用，已知化学反应会强烈改变对流的起始。这些是未来研究的课题。
	
	尽管存在这些局限性，从 YUCT 循环中代数推导出的临界协调效率 \(K_{\mathrm{eff}}^{\mathrm{crit}} \approx 1.837\) 在框架内是精确的。大自然是否实际使用这个值可以通过实验来决定。
	
	\section{结论}
	
	我们已经证明 YUCT 框架为理解 B\'{e}nard 对流的起始提供了定量的、可检验的解释。临界 Rayleigh 数用流体的协调效率表达，而协调效率本身又与可测量的输运系数相联系。预测的标度关系 \(\mathrm{Ra_c} \propto \eta/\lambda\) 可以用现有的实验数据验证。这项工作表明 YUCT 不仅仅是热力学的一种哲学重述，而是一个能够产生可证伪的数值预测的框架。
	
	\section*{致谢}
	作者感谢 YUCT 研讨会参与者的批判性讨论。
	
	\begin{thebibliography}{99}
		\bibitem{Yakushev2026} A.~V.~Yakushev, \emph{Yakushev Unified Coordination Theory (YUCT)}, Zenodo, DOI:10.5281/zenodo.18444599 (2026).
		\bibitem{AppendixY} A.~V.~Yakushev, ``Appendix Y: Mathematical Foundations of YUCT,'' in \emph{YUCT}, Zenodo (2026).
		\bibitem{AppendixL} A.~V.~Yakushev, ``Appendix L: Fractal Coordination Error Scaling,'' in \emph{YUCT}, Zenodo (2026).
		\bibitem{AppendixFerra} A.~V.~Yakushev, ``Appendix Ferra1.2: Universal Coordination Constants for Ordered and Disordered Phases,'' in \emph{YUCT}, Zenodo (2026).
		\bibitem{AppendixAF} A.~V.~Yakushev, ``Appendix AF: The Algebraic Framework of Reality in YUCT,'' in \emph{YUCT}, Zenodo (2026).
		\bibitem{Prigogine1977} I.~Prigogine and G.~Nicolis, \emph{Self‑Organization in Nonequilibrium Systems}. Wiley (1977).
		\bibitem{Glansdorff1971} P.~Glansdorff and I.~Prigogine, \emph{Thermodynamic Theory of Structure, Stability and Fluctuations}. Wiley (1971).
		\bibitem{Chandrasekhar1961} S.~Chandrasekhar, \emph{Hydrodynamic and Hydromagnetic Stability}. Oxford (1961).
		\bibitem{Cross1993} M.~C.~Cross and P.~C.~Hohenberg, ``Pattern formation outside of equilibrium,'' \emph{Rev. Mod. Phys.} \textbf{65}, 851 (1993).
	\end{thebibliography}
	
\end{document}